\section{A one-sided test on sparse twists and sparse scales}
\label{sec:sparse}
A separate consequence uses the proof's zero-detection argument before its
final GRH theorem. The published Lean criterion takes an eventual inequality
for every auxiliary prime in one residue class and every sufficiently large
real scale. Both sets of inputs can be thinned. The conclusion below is a
reduction proved here on paper. Sandi Setiawan already gives the fourth-power
mesh and a chord interpolation argument for the ordinary zeta prime sum
(\emph{The Primacohedron Trace Formulation of the Riemann Hypothesis},
22 November 2026, Theorems 3.2 and 4.1,
\href{https://www.academia.edu/181628231}{author's manuscript}).
Here the additional reduction selects a thin family of quadratic twists
using the repository's quantitative zero-density bound. We make no claim
of priority for the combined criterion.

Fix an integer $Q>0$ divisible by $8$ and a reduced residue class
$a\pmod Q$ with $a\equiv7\pmod8$. For primes $d\equiv a\pmod Q$, put
\[
 w_d(n)=\mathbf1_{n\equiv1\bmod Q}\,\Lambda(n)
       \bigl(1-(n/d)\bigr),\qquad
 R_d(x)=\sum_{1\le n\le x} w_d(n)(1/n-1/x)
       -\frac{\log x}{\varphi(Q)}-c_{d,Q}.
\]
Here $(n/d)$ is the Legendre symbol and $c_{d,Q}$ is the unique constant
for which $R_d(x)\to0$. The repository gives its explicit logarithmic
derivative formula in \texttt{New.logDerivMean} and identifies it in
\texttt{New.primeRemainder}. All $w_d(n)$ are nonnegative.

\begin{theorem}[Sparse one-sided zero detection]
Fix $1/2<\beta<1$. Let $S$ be a set of the above primes and
$S(X)=\#\{d\in S:d\le X\}$. Suppose
\[
 \limsup_{X\to\infty}\frac{\log(1+S(X))}{\log X}
 \ \ge\ 1-\frac{\beta-1/2}{18}.
 \tag{S}
\]
For each $d\in S$, suppose there is an increasing unbounded sequence
$x_j\ge1$, allowed to depend on $d$, with
\[
 x_{j+1}-x_j=O_d(x_j^{(1+\beta)/2}),\qquad
 R_d(x_j)\le C_d x_j^{\beta-1}\quad(j\text{ sufficiently large}),
 \tag{E}
\]
where $C_d>0$. For every Dirichlet character $\psi\bmod Q$,
$L(s,\psi)$ is nonzero in $\Re s>\beta$, with the principal pole excepted.
\end{theorem}

\begin{proof}
First recover all intervening scales. Write $t=1/x$. On any compact interval
$[u,v]\subset(0,\infty)$ the function
\[
 F(t)=\sum_{n\ge1}w_d(n)\,(1/n-t)_+
\]
is a finite sum of convex functions, hence lies below its chord. We have
$R_d(1/t)=F(t)+(\log t)/\varphi(Q)-c_{d,Q}$. Write $I f$ for the affine
interpolant from $u,v$. Since $t\mapsto\log t+t^2/(2u^2)$ is convex on
$[u,v]$,
\[
 0\le\log t-I(\log)(t)
 \le\frac{(t-u)(v-t)}{2u^2}
 \le\frac{(v-u)^2}{8u^2}.
\]
Take $u=1/x_{j+1}, v=1/x_j$. The ratio $(v-u)/u$ equals
$(x_{j+1}-x_j)/x_j$, so the interpolation loss is
$O_d(x_j^{\beta-1})$. Moreover $x_{j+1}/x_j\to1$. Applying (E) at
the endpoints gives $R_d(x)\le K_d x^{\beta-1}$ throughout the tail for
some $K_d>0$.

Suppose $L(\rho,\psi)=0$ with $\sigma=\Re\rho>\beta$. At $\rho=1$
the classical Dirichlet nonvanishing theorem applies to nonprincipal
characters. The ordinary zero-free region $\Re s\ge1$ therefore leaves
$\sigma<1$. Choose $\gamma$ with $\beta<\gamma<\sigma$ and
set $H=|\Im\rho|$. Call an admissible prime $d$ bad if at least one
$L(s,\xi(\cdot/d))$, $\xi\bmod Q$, has a zero in
$\Re s\ge\gamma$, $|\Im s|\le H$. The upstream lemma
\texttt{New.zero\_density\_quadTwist}, with
$\delta=(\gamma-1/2)/9$, gives
\[
 \#\{d\le X:d\text{ bad}\}
   = O_{Q,a,\gamma,H}(X^{1-(\gamma-1/2)/18}).
\]
This exponent is strictly smaller than the right side of (S).
Thus some $d\in S$ lies outside the bad set. Fix this $d$.
The real-variable estimate already proved implies
$R_d(x)\le x^{\gamma-1}$ eventually, since $K_d x^{\beta-\gamma}\to0$.

For clarity, here is where one good twist suffices. The meromorphic
transform of $v\mapsto R_d(e^v)$ in \texttt{New.isPolar\_primeRemainder}
has residue at $s\ne1$ equal to
\[
 \frac{m(s)-m_d(s)}{\varphi(Q)s(1-s)},
\]
where $m$ and $m_d$ sum the zero multiplicities of the pole-removed
untwisted and twisted functions over all $\xi\bmod Q$. The transform is
regular at $1$. In the box, $m_d=0$. On its real segment
$\gamma<s<1$ the residues are nonnegative; for real $s\ge1$ they
vanish. The exponential upper bound $R_d(e^v)\le e^{(\gamma-1)v}$
and these real signs are precisely the hypotheses of the repository's
one-sided Landau lemma, \texttt{IsPolar.landau}. It forces the residue
at $\rho$ to vanish. Then $m(\rho)=0$, contrary to the stipulated zero of
$L(s,\psi)$. This is the fixed-$d$ part of the proof of
\texttt{New.LFunction\_ne\_zero\_of\_primeRemainder\_le}.
\end{proof}

\paragraph{What a counterexample would force.}
If a zero of real part $\sigma>\beta$ existed, let $T_\beta(X)$ count
admissible primes $d\le X$ for which
$\limsup_{j\to\infty}j^{4(1-\beta)}R_d(j^4)<+\infty$.
As $j^4$ meets (E)'s gap restriction, the proof actually gives
\[
 T_\beta(X)=O_{Q,a,\rho,\gamma}
 \bigl(X^{1-(\gamma-1/2)/18}\bigr)
 \qquad(\beta<\gamma<\sigma).
\]
Indeed, every such $d$ would have to be bad in that single box.
Thus the discrepancy at fourth powers would have arbitrarily large
positive normalized excursions for every admitted prime outside
an exceptional set of power-saving size.

\paragraph{An explicit sampling schedule.}
Enumerate the primes in the prescribed class as $p_1<p_2<\cdots$.
Choose $d_k=p_{\lfloor k\log^2(k+2)\rfloor}$ for all sufficiently large
$k$. The prime number theorem in progressions gives
$d_k\sim\varphi(Q)k\log^3 k$. Their reciprocals have a convergent sum,
and their counting exponent in (S) is $1$. Take $x_j=j^4$.
The gap is $O(x_j^{3/4})$. If for every selected $d_k$ and every
$\varepsilon>0$ the sampled values satisfy
\[
 R_{d_k}(j^4)\le C_{k,\varepsilon}\,j^{-2+4\varepsilon}
 \quad\text{for all sufficiently large }j,
\]
apply the theorem with $\beta=1/2+\varepsilon$ for arbitrarily small
positive $\varepsilon$. This excludes all zeros to the right of $1/2$
modulo $Q$ apart from the principal pole. For arbitrary original modulus
$q$, take $Q=18480q$, $a=Q/2-1$; this class meets the conditions because
$Q$ is divisible by $16$ and $\gcd(Q/2-1,Q)=1$. Changing the level adds
finite Euler factors, so the nonvanishing descends to $q$. For these
parameters the upstream one-sided estimate already implies the sampled
inequalities. Their interest is the weaker input an alternative proof
would have to establish. Bounds for finitely many twists or finitely many
scales do not give the required asymptotic estimates.
